\documentclass[11pt,a4paper]{article}
\usepackage[utf8]{inputenc}
\usepackage[T1]{fontenc}
\usepackage[german]{babel}
\usepackage{geometry}
\usepackage{titlesec}
\usepackage{fancyhdr}
\usepackage{enumitem}

\geometry{a4paper, left=25mm, right=25mm, top=25mm, bottom=25mm}


\titleformat{\section}{\Large\bfseries}{\thesection}{1em}{}
\titleformat{\subsection}{\large\bfseries}{\thesubsection}{1em}{}

\pagestyle{fancy}
\fancyhf{}
\renewcommand{\headrulewidth}{0pt}

\begin{document}

\begin{center}
    \LARGE\textbf{Association Statutes}
    \\
    \small{of the}
    \\
    \large{Math Olympiad Club Zurich}
    
    \vspace{0.5cm}
    \textbf{Updated on 22.10.2025}
\end{center}

\vspace{1cm}

\section*{1. Name and Registered Office}

An association within the meaning of Articles 60 et seq. of the Swiss Civil Code is established under the name:
\begin{center}
    Math Olympiad Club Zurich
\end{center}
with its registered office in Zurich. It is politically independent and denominationally neutral.

\section*{2. Purpose and Objectives}

The association's purposes are:

\begin{enumerate}
    \item to motivate members in solving competitive mathematical problems, puzzles;
    \item to promote mathematical projects and project work by students in Zurich, particularly those of students from ETH Zurich and the University of Zurich.
\end{enumerate}

\section*{3. Resources}

For the pursuit of its purpose, the association has the following resources:

\begin{enumerate}
    \item subsidies;
    \item income from its own events;
    \item donations and contributions of any kind.
\end{enumerate}

No membership fees are collected.

\section*{4. Members}

Natural persons and legal entities who are interested in the purpose of the association or who are supported by the association may become members.
\\\\
Active members with voting rights are natural persons who use the association's offerings.
\\\\
Passive members with voting rights may be natural persons or legal entities who support the association ideologically and financially.
\\\\
The committee may confer the title of honorary member upon persons who have made a special commitment to the association. They do not have voting rights.
\\\\
Membership in the association can be applied for at any time. Applications for admission must be submitted to the committee; the committee has the final decision on admission.
\section*{5. Termination of Membership}

Membership ends:

\begin{enumerate}
    \item for natural persons: by resignation, exclusion, or death;
    \item for legal entities: by resignation, exclusion, or dissolution of the legal entity.
\end{enumerate}

\section*{6. Resignation and Exclusion}

Resignation from the association is possible at any time by notifying the committee. The resignation letter must be submitted in writing to the committee at least 2 weeks before the intended date.
\\\\
A member may be excluded at any time by the committee without stating reasons.

\section*{7. Bodies of the Association}

The bodies of the association are:

\begin{enumerate}
    \item the General Assembly of members;
    \item the Committee.
\end{enumerate}
\section*{8. The General Assembly}

The supreme body of the association is the General Assembly. An ordinary General Assembly takes place twice per year, at the beginning of each new university semester. The relevant periods are: September-November and January-April.
\\\\
Members are convened to the General Assembly at least 10 days in advance in writing. Convening via email and WhatsApp is valid.
\\\\
Proposals from members for additional agenda items for the General Assembly must be submitted in writing to the committee, with justification, at least one week before the assembly.
\\\\
The committee or 1/5 of the members may at any time request the convening of an extraordinary General Assembly by stating its purpose. The assembly must be held no later than 4 weeks after receipt of the request.
\\\\
The General Assembly has the following mandatory tasks and competencies:

\begin{enumerate}
    \item Approval of the minutes of the previous General Assembly;
    \item Approval of the committee's annual report;
    \item Review of the annual accounts;
    \item Granting of discharge to the committee;
    \item Election of the President and other committee members;
    \item Decision on the program of activities;
    \item Decision on proposals from the committee and members;
    \item Amendment of the statutes;
    \item Decision on the dissolution of the association and the use of the liquidation proceeds.
\end{enumerate}
Any duly convened General Assembly has the power to deliberate, regardless of the number of members present.
\\\\
Members make decisions by a simple majority of the votes cast. In the event of a tie, the President has the casting vote.
\\\\
Statutory amendments require the approval of a 2/3 majority of the voting members present.
\\\\
Minutes must be kept of at least the decisions taken.

\section*{9. The Committee}

The committee is composed of at least one person. The term of office extends until the next General Assembly at the beginning of the semester. The committee manages the day-to-day affairs and represents the association vis-à-vis third parties.
\\\\
It issues regulations, can create working groups, and can employ (according to labor law) or mandate persons against appropriate compensation to achieve the association's objectives.
\\\\
The committee has all competencies that are not assigned by law or by these statutes to another body.
\\\\
The committee constitutes itself.
\\\\
The committee meets as often as business requires. Any committee member may request the convening of a meeting by stating the reasons. If no committee member requests an oral deliberation, decision-making by circular resolution (including via email and WhatsApp) is valid.
\\\\
The committee operates essentially on a voluntary and unpaid basis. It is entitled to reimbursement of actual expenses. Appropriate compensation may be paid for special services by individual committee members.

\section*{10. Signatory Authority}

The association is legally bound by the joint signature of the President together with another committee member.

\section*{11. Liability}

Only the assets of the association are liable for the debts of the association. Any personal liability of the members is excluded.

\section*{12. Data Protection}

The association only collects from members the personal data necessary to fulfill its purpose. The committee ensures data security appropriate to the risk.
\\\\
Member data is not disclosed to other members, unless required by a legal provision. Furthermore, data is only disclosed to third parties within the framework of lawful processing and if it is legally required or ordered by an authority.
\\\\
The processing of member data is otherwise carried out in accordance with the provisions of Swiss data protection legislation and the privacy policy on the association's website.

\section*{13. Dissolution of the Association}

The dissolution of the association can be decided by a resolution of an ordinary or extraordinary General Assembly with a 2/3 majority of the members present.
\\\\
In the event of the dissolution of the association, the assets of the association shall devolve to a tax-exempt organization in Switzerland which pursues the same or a similar purpose. The distribution of the association's assets among the members is excluded.

\section*{14. Entry into Force}

These statutes were adopted at the constitutive assembly on 22.10.2025 and entered into force on that date.

\vspace{2cm}

\begin{flushright}
Place: ETHZ HG, Zurich
\\
Date: 22.10.2025

\vspace{1cm}

Antoine du Fresne von Hohenesche (President)

\end{flushright}
\end{document}